% !PRINTED 175
\weberpage{175}
Dr\"uckt man in diesen Gleichungen $\omega'$ nach (9) durch $\omega$ aus, so erh\"alt man

\[
\begin{aligned}
1 &= \alpha'(\alpha+\beta\omega) + \beta'(\gamma+\delta\omega),\\
\omega &= \gamma'(\alpha+\beta\omega) + \delta'(\gamma+\delta\omega),
\end{aligned}
\]

und da $\omega$ nicht reell ist, so zerf\"allt jede dieser Gleichungen in zwei andere:

\[
\begin{aligned}
\alpha\alpha' + \gamma\beta' &= 1, & \beta\alpha' + \delta\beta' &= 0,\\
\alpha\gamma' + \gamma\delta' &= 0, & \beta\gamma' + \delta\delta' &= 1.
\end{aligned}
\]

Mithin ist

\[
(\alpha\delta - \beta\gamma)(\alpha'\delta' - \beta'\gamma') = 1,
\]

und daher, da $\omega,\omega'$ beide einen positiven imagin\"aren Teil haben m\"ussen, also nach (9) die Determinante $\alpha\delta-\beta\gamma$ positiv sein mu\ss{}:

\[
\alpha\delta - \beta\gamma = 1,
\]
\[
\alpha = \delta', \quad \delta = \alpha', \quad \gamma = -\gamma', \quad \beta = -\beta'.
\]

1. \gesperrt{Es h\"angen also $\omega, \omega'$ durch eine lineare Transformation, und zwar [nach (6), (8)] durch eine der ersten Klasse (\S\ 36) miteinander zusammen. Da\ss{} auch umgekehrt zwei solche Werte $\omega, \omega'$ denselben Wert $\varkappa^2$ ergeben, ist bereits oben nachgewiesen.}

Ein \"ahnlicher Satz ergibt sich als unmittelbare Konsequenz hieraus, f\"ur die Invariante $j(\omega)$ (\S\ 46).

2. \gesperrt{Die Gleichung:}

\weq{10}{j(\omega') = j(\omega)}

\gesperrt{ist dann und} \textbf{nur} \gesperrt{dann befriedigt, wenn $\omega'$ mit $\omega$ durch irgend eine lineare Transformation}

\[
\begin{pmatrix} \alpha, & \beta \\ \gamma, & \delta \end{pmatrix}
\]

\gesperrt{zusammenh\"angt, wenn also}

\weq{11}{\omega' = \frac{\gamma+\delta\omega}{\alpha+\beta\omega}, \quad \alpha\delta-\beta\gamma = 1}

\gesperrt{ist.}

Denn bezeichnen wir die zu $\omega'$ geh\"origen Werte von $\varkappa^2,\varkappa'^2$ mit $\lambda^2,\lambda'^2$, so kann die Gleichung (10) so geschrieben werden:

\[
\frac{(1-\lambda^2\lambda'^2)^3}{\lambda^4\lambda'^4} = \frac{(1-\varkappa^2\varkappa'^2)^3}{\varkappa^4\varkappa'^4},
\]

und ist in bezug auf $\lambda^2$ eine Gleichung sechsten Grades, deren Wurzeln

\[
\varkappa^2, \quad \varkappa'^2, \quad \frac{1}{\varkappa^2}, \quad \frac{1}{\varkappa'^2}, \quad -\frac{\varkappa'^2}{\varkappa^2}, \quad -\frac{\varkappa^2}{\varkappa'^2}
\]
%CONT>>
